% =============================================================================
% Lecture 0: Course Overview
% Physical AI: Machine Learning Systems That Sense and Act (Volume IV)
% =============================================================================
% The first day of class: what physical AI is, why delay and evidence are paid
% in clearance, what the machine is made of, and how the course runs. Lecture 0
% owns no concepts. It previews the elementary stopping budget and the stack;
% Lecture 1 owns both (course plan, "What Lecture 0 leaves to Lecture 1").
%
% Running case: the warehouse mobile manipulator on its aisle run (machine card,
% The Causal Boundary, table 2), shown as a real machine of the same class by
% a photo of an ER-FLEX mobile manipulator (MiR250 base, UR16e arm).
%
% Build: make deck V=4 D=00_course_overview   (figures: images/make_figs.py)
% Numbers: images/calcs.py prints every derived value on these slides.
% Plan: Volume IV course plan, "Lecture 0: Course overview".
% =============================================================================

\documentclass[aspectratio=169, 12pt]{beamer}
\usepackage{../../assets/beamerthememlsys}
\mlsyssetup{
  volume={Volume IV},
  chapter={Lecture 0},
  chaptertitle={Course Overview},
  logo={../../assets/img/logo-mlsysbook.png},
  instlogo={../../assets/img/logo-harvard.png},
}
\usepackage[T1]{fontenc}
\usepackage[scaled=0.9]{helvet}
\usepackage{courier}
\renewcommand{\familydefault}{\sfdefault}
\usepackage{booktabs,amsmath}
\graphicspath{{images/}}
\setcounter{mlsystotalsections}{6}

\title{Physical AI}
\author{Vijay Janapa Reddi}
\institute{Harvard University}
\date{}

% A large answer choice for predict and concept-question slides
\newcommand{\choice}[2]{\textbf{\textcolor{mlsysaccent}{#1}}\enspace #2}
% Volume IV roles (volumes/vol4/palette.md)
\definecolor{delaycol}{HTML}{D97706}
\definecolor{plantcol}{HTML}{334155}
\definecolor{nervecol}{HTML}{059669}
\definecolor{conseq}{HTML}{DC2626}
\definecolor{brainnavy}{HTML}{1F407A}

\begin{document}

% =============================================================================
% TITLE
% =============================================================================
\mlsysnote{[1 min] Welcome. Read the thesis aloud once. The photo is a real mobile manipulator, a wheeled base carrying an arm; it is the class of machine the course follows all term, and it returns several times today. Today previews the course through one question, how far a robot travels while it decides, and gives students two pictures to keep: the five-level stack and the stopping-budget bar. Lecture 1 owns every idea previewed today. The notes total 70 minutes; cutting every [OPTIONAL] note leaves a lecture of about 58 minutes.}
\mlsystitle{Physical AI}{Once a learned model's output moves mass, the time and evidence behind each action are paid in clearance, so an independent path must be able to refuse the action before that clearance runs out.}{erflex-mobile-manipulator.jpg}

% =============================================================================
% SECTION: ONE FRAME
% =============================================================================
\mlsysnote{[0.5 min] NARRATE: open with a question every student can answer with a guess: how far does a robot move while it thinks? FLEX: [CORE].}
\mlsysdivider{How far does a robot travel while it decides?}{The Causal Boundary, \S1.1 and \S1.4.3}
\section{One Frame}

\begin{frame}{Predict: how far does the base move during one frame?}
  \mlsysactivity{Predict}{2.5 min}
  \begin{columns}[c]
    \begin{column}{0.48\textwidth}
      \centering\safeimg[width=\textwidth,height=4.8cm,keepaspectratio]{erflex-mobile-manipulator.jpg}
    \end{column}
    \begin{column}{0.48\textwidth}
      \small A base like this drives at 1.5 m/s. A person steps out at the rack end. How far does the base travel before perception hands on one camera frame?\\[0.2cm]
      \large
      \choice{A}{1 mm}\qquad\choice{B}{10 mm}\\[0.15cm]
      \alt<2->{\colorbox{mlsysaccent!15}{\choice{C}{80 mm}}}{\choice{C}{80 mm}}\qquad\choice{D}{800 mm}
    \end{column}
  \end{columns}
  \mlsyscredit{Photo: ER-FLEX mobile manipulator (MiR250 base, UR16e arm), Auledas, CC BY 4.0, Wikimedia Commons. Speed: the course's running machine, The Causal Boundary, table 2 (illustrative).}
  \note{[2.5 min] NARRATE: the photo shows a real mobile manipulator, an industrial arm on an autonomous wheeled base. The course's running machine is a warehouse robot of this class; its numbers come from the book's machine card rather than from this product. Ask for a vote with fingers, one to four, then one minute with a neighbor, then vote again. Step 2 reveals C: about 80 mm, exactly 77.4 mm, derived next. ENGAGE: ask a student who chose A what they pictured. Expected: the inference time of a fast model, a few milliseconds. WARN: students picture only the neural network. The frame ages while the camera exposes and reads out, while bytes cross the bus, and while the image is processed; D confuses one frame with the whole stop, which comes later. FLEX: [CORE].}
\end{frame}

\begin{frame}{The base covers 77.4 mm before one frame is dispatched}
  \centering
  \mlsyssteps[width=0.95\textwidth,height=5.6cm,keepaspectratio]{aisle_step}{2}
  \mlsyscredit{The Causal Boundary, \S1.4.3; Sensor Perception, Napkin Math 8.2. Running-machine values (illustrative and derived); computed in images/calcs.py.}
  \note{[2 min] LINK: the vote. NARRATE: Step 1: the aisle seen from above. The base drives right at 1.5 m/s toward the end of a rack. A person may step out there. The red line marks the clear distance, 1,100 mm, the room the base has to stop in. Step 2: the camera frame that first shows the person is 51.6 ms old when perception dispatches it. At 1.5 m/s that is 1.5 times 51.6, or 77.4 mm, traveled before any later stage has seen the person. ENGAGE: ask what a model twice as fast would save. Expected: only its own share of the 51.6 ms; the next slide shows that share. WARN: the 77.4 mm is not a stop; it is only the age of one frame. FLEX: [CORE]. Lecture 8 owns observation age.}
\end{frame}

\begin{frame}{A frame is 24.6 ms old before any computation starts}
  \centering
  \safeimg[width=0.95\textwidth,height=3.4cm,keepaspectratio]{frame_age.pdf}\\[0.3cm]
  {\small Camera exposure and readout take as long as the learned vision network.}
  \mlsyscredit{Sensor Perception, Napkin Math 8.2 (readout datasheet; other stages chosen or illustrative); computed in images/calcs.py.}
  \note{[1.5 min] NARRATE: the 51.6 ms by stage. The first 24.6 ms come from the camera itself: half of a 16 ms exposure, measured from its midpoint, plus 16.6 ms to read the rows out. Transfer, queueing, and image processing add 4.5 ms. The learned backbone, an INT8 vision network, takes 22 ms, and dispatch 0.5 ms. ENGAGE: ask which stage a better model can shorten. Expected: only the backbone, 22 of 51.6 ms. WARN: the stage values are budgets the rest of the machine plans against; none is a measurement of a robot. FLEX: [OPTIONAL]; if cut, say ``half the frame's age is the camera'' on the previous slide. Lecture 8 builds this waterfall stage by stage.}
\end{frame}

\begin{frame}{A software agent can wait and roll back; a moving base cannot}
  \begin{columns}[c]
    \begin{column}{0.62\textwidth}
      \small
      \begin{tabular}{@{}P{2.2cm}P{2.8cm}P{2.9cm}@{}}
        \toprule
        & \textbf{Behind the glass} & \textbf{In the aisle} \\
        \midrule
        \textbf{Time} & waits 200 ms or 30 s for the model & keeps moving at 1.5 m/s \\[0.1cm]
        \textbf{Mistakes} & roll back the transaction & \textcolor{conseq}{\textbf{cannot be recalled}} \\[0.1cm]
        \textbf{Interface} & strings checked before execution & current, torque, contact \\
        \bottomrule
      \end{tabular}
    \end{column}
    \begin{column}{0.35\textwidth}
      \centering\safeimg[width=\textwidth,height=4.4cm,keepaspectratio]{erflex-mobile-manipulator.jpg}
    \end{column}
  \end{columns}
  \mlsyscredit{The Causal Boundary, \S1.1 and table 1. Photo: Auledas, CC BY 4.0, Wikimedia Commons.}
  \note{[2 min] LINK: the 77.4 mm happened because the world did not wait. NARRATE: a software agent working in a controlled environment, ``behind the glass,'' keeps three privileges. Its environment can wait while the model computes, 200 ms or 30 s. A failed action can be rolled back. Its actions are strings that a checker can validate before they run. A machine that moves mass loses all three: physical time keeps running, work delivered to the world cannot be recalled, and its outputs become current and torque. ENGAGE: ask for an example of an action that software cannot undo. Expected: a dent, a fracture, heat in a winding, a person struck. WARN: behind-the-glass agents also lose the privileges when a tool call reaches the outside world; the table describes the controlled environment. FLEX: [CORE]. Lecture 1 owns the privileges lost to embodiment.}
\end{frame}

\begin{frame}{Driving question: how far does a robot travel while it decides?}
  \begin{columns}[c]
    \begin{column}{0.68\textwidth}
      \small
      Today you will be able to
      \begin{itemize}\setlength{\itemsep}{0.12cm}
        \item tell physical AI from advisory learning and classical control;
        \item compute how much stopping clearance a delay spends;
        \item name the five levels of the machine and its two boundaries;
        \item explain why a clean test record is weak evidence, and name the four laws.
      \end{itemize}
    \end{column}
    \begin{column}{0.29\textwidth}
      \centering
      \safeimg[width=\textwidth,height=1.05cm,keepaspectratio]{thumb_aisle_step2.pdf}\\[0.12cm]
      \safeimg[width=\textwidth,height=1.05cm,keepaspectratio]{thumb_stop_bar_step3.pdf}\\[0.12cm]
      \safeimg[width=\textwidth,height=1.05cm,keepaspectratio]{thumb_stack.pdf}\\[0.12cm]
      \safeimg[width=\textwidth,height=1.05cm,keepaspectratio]{thumb_clean_runs_step4.pdf}
    \end{column}
  \end{columns}
  \note{[1.5 min] NARRATE: read the driving question; the last segment answers it. Each goal is one segment, and each thumbnail is the picture that teaches it: the aisle, the stopping-budget bar, the stack, and the clean-run ladder. The four laws are the course's spine: the first law, irreversibility, is physical; the other three are design principles. FLEX: [CORE].}
\end{frame}

% =============================================================================
% SECTION: ORIGINS
% =============================================================================
\mlsysnote{[0.5 min] NARRATE: before the arithmetic, what counts as physical AI and where it came from. FLEX: [CORE].}
\mlsysdivider{What makes a machine physical AI, and where did it come from?}{The Causal Boundary, \S1.2 and \S1.6}
\section{Origins}

\begin{frame}{Physical AI needs all three: a learned model, feedback, authority}
  \centering
  \safeimg[width=0.8\textwidth,height=5.8cm,keepaspectratio]{scope_venn.pdf}
  \mlsyscredit{After The Causal Boundary, \S1.6, figure 6.}
  \note{[2 min] NARRATE: the scope test has three criteria. A learned model: behavior that comes from training data, so its correctness has to be shown by testing. Consequential physical feedback: actions change the world, and the world changes what the sensors see next. Delegated authority: outputs reach the actuators without a person approving each one. Each overlap of two circles is something else. A learned model inspected by a person before it acts is advisory. Learned outputs that run unapproved in software are digital autonomy. Physical feedback without learning is classical control. ENGAGE: ask whether adding a certified safety brake removes delegated authority. Expected: no; learned proposals still drive all motion up to the brake. WARN: students think the model must write the motor registers itself; a learned detector that triggers braking already shapes actuation. FLEX: [CORE]. Lecture 1 owns the scope test.}
\end{frame}

\begin{frame}{Which of these is physical AI?}
  \mlsysactivity{Vote}{3 min}
  \begin{columns}[T]
    \begin{column}{0.32\textwidth}
      \centering
      \begin{mlsyscard}{mlsysaccent}\small\centering
        A learned model labels warehouse photos; a person reads the labels
      \end{mlsyscard}
      \vspace{0.1cm}{\Large\choice{A}{}}
    \end{column}
    \begin{column}{0.32\textwidth}
      \centering\safeimg[width=\textwidth,height=2.3cm,keepaspectratio]{kuka-spot-welding-line.jpg}\\
      {\footnotesize Fixed arms spot-weld car bodies}\\[0.05cm]{\Large\choice{B}{}}
    \end{column}
    \begin{column}{0.32\textwidth}
      \centering\safeimg[width=\textwidth,height=2.3cm,keepaspectratio]{erflex-mobile-manipulator.jpg}\\
      {\footnotesize A learned policy drives base and arm}\\[0.05cm]{\Large\alt<2->{\colorbox{mlsysaccent!15}{\choice{C}{}}}{\choice{C}{}}}
    \end{column}
  \end{columns}
  \vspace{0.05cm}
  \uncover<2->{\centering\small A is advisory. B is classical control. Only C meets all three criteria.\par}
  \mlsyscredit{Photos: BMW Werk Leipzig (KUKA robots), CC BY-SA 2.0 de; Auledas, CC BY 4.0; Wikimedia Commons. The Causal Boundary, \S1.6.}
  \note{[3 min] NARRATE: vote, one minute with a neighbor, vote again. Step 2: only C. A has a learned model but no delegated authority; a person decides. B, KUKA arms on a BMW body line, closes a physical loop with no learned model in the motion; the arms replay taught motion under classical controllers. C uses a learned policy whose outputs reach the drives without approval of each command. ENGAGE: ask what would move B into scope. Expected: a learned vision model that changes where each weld goes, with no person approving each change. WARN: ``it is a robot'' is not the test; B is a robot and is not physical AI. FLEX: [CORE].}
\end{frame}

\begin{frame}{The learned model moved step by step toward the actuator}
  \begin{columns}[b]
    \begin{column}{0.30\textwidth}
      \centering\safeimg[width=\textwidth,height=2.3cm,keepaspectratio]{unimate-1967.jpg}\\
      {\footnotesize Unimate: replays stored steps}
    \end{column}
    \begin{column}{0.30\textwidth}
      \centering\uncover<2->{\safeimg[width=\textwidth,height=2.3cm,keepaspectratio]{shakey-1969.jpg}\\
      {\footnotesize Shakey: plans, then acts}}
    \end{column}
    \begin{column}{0.30\textwidth}
      \centering\uncover<4->{\safeimg[width=\textwidth,height=2.3cm,keepaspectratio]{erflex-mobile-manipulator.jpg}\\
      {\footnotesize Learned policies propose}}
    \end{column}
  \end{columns}
  \vspace{0.15cm}
  \centering
  \mlsyssteps[width=0.92\textwidth,height=2.2cm,keepaspectratio]{learned_band_step}{4}
  \mlsyscredit{Photos (Wikimedia Commons): Unimate, 1967, Frank Q. Brown, LA Times, CC BY 4.0; Shakey, The wub, CC BY-SA 4.0; Auledas, CC BY 4.0. Band: The Causal Boundary, \S1.2, fig.\ 1.}
  \note{[2.5 min] NARRATE: the band shows where the learned output lands in typical deployed machines. Step 1: the Unimate arm began work at General Motors in 1961 by replaying commands stored on a magnetic drum; nothing was learned, and people were kept out of its reach. The photo is a Unimate at a 1967 demonstration. Step 2: Shakey, at SRI from 1966 to 1972, perceived and planned, but the world was assumed to wait while it thought. Step 3: from the mid-1980s learned perception fed hand-written planners (ALVINN steered a van in 1989; Stanley won the 2005 Grand Challenge). Step 4: from about 2016 networks emitted motor commands, and from 2023 pretrained vision-language-action models propose motion. ENGAGE: ask what had to move inside the machine at each step. Expected: the protection that separation or a person used to supply. WARN: the 1967 photo shows a public demonstration of pouring coffee rather than a factory cell. FLEX: [CORE]. The reading (\S1.2) tells the history in full.}
\end{frame}

\begin{frame}{Predict: how much less compute do robot policies use?}
  \mlsysactivity{Predict}{2.5 min}
  \centering
  \mlsyssteps[width=0.95\textwidth,height=5.0cm,keepaspectratio]{epoch_compute_step}{2}
  \only<1|handout:0>{\par\small A 2024 robot policy against the largest language model of 2024: 3$\times$, 30$\times$, 300$\times$, or 3,000$\times$ less?}%
  \only<2-|handout:1>{\par\small Only 25 robotics models in the database carry a compute estimate.}
  \mlsyscredit{Epoch AI, Data on AI Models (2026), via MLSysIM Literature.ModelsByDomain; The Causal Boundary, \S1.2, figure 2. Gaps computed in images/calcs.py.}
  \note{[2.5 min] NARRATE: Step 1: each dot is a model with an Epoch AI estimate of its training compute, on a log axis. Gray are language models, green are robotics models. Take the vote on the question below the plot. Step 2: OpenVLA, a 2024 open vision-language-action model, trained on about 350 times less compute than the largest 2024 language model (calcs.py prints 345); GR00T N1, a 2025 open robot foundation model, about 7,000 times less. ENGAGE: ask why. Expected: robot data must be collected by machines acting in the world, and the model must run on the robot within its power and time budget. WARN: compute estimates are Epoch's; only 25 robotics models carry one, so the green points are a thin record. FLEX: [CORE]. The reading (\S1.2) states the two constraints.}
\end{frame}

% =============================================================================
% SECTION: THE BUDGET
% =============================================================================
\mlsysnote{[0.5 min] NARRATE: now the arithmetic. Where does a command become work, and how much clearance does delay spend? FLEX: [CORE].}
\mlsysdivider{Where does a command become physical work, and what does delay cost there?}{The Causal Boundary, \S1.4}
\section{The Budget}

\begin{frame}{A command becomes physical work when current enters a winding}
  \begin{columns}[c]
    \begin{column}{0.62\textwidth}
      \centering\mlsyssteps[width=\textwidth,height=5.4cm,keepaspectratio]{chain_step}{4}
    \end{column}
    \begin{column}{0.35\textwidth}
      \centering\safeimg[width=\textwidth,height=3.6cm,keepaspectratio]{bldc-stator-winding.jpg}\\[0.05cm]
      {\footnotesize Copper windings of a brushless DC motor's stator}
    \end{column}
  \end{columns}
  \mlsyscredit{After The Causal Boundary, \S1.4.1 and Definition 1.2. Photo: Medvedev, CC BY-SA 3.0, Wikimedia Commons.}
  \note{[2.5 min] NARRATE: Step 1: an admitted command is a setpoint, a number written into a register on a microcontroller. Step 2: the register sets how long power switches stay on, and current flows in the stator windings; point to the copper coils in the photo. Step 3: that current makes a magnetic field that turns the rotor, torque on the shaft. Step 4: torque accelerates the base, the arm, and the payload. The dashed line is the causal boundary: the interface through which a digital command produces physical forces, currents, or flows. Above it a value can be rewritten before the next clock cycle; below it, work delivered cannot be recalled, and a correction is another motion with its own time and distance. ENGAGE: ask whether killing the policy's process stops the base. Expected: no; the base keeps its momentum until something commands a stop. FLEX: [CORE]. Lecture 1 owns the causal boundary.}
\end{frame}

\begin{frame}{The running machine: a 368 kg base and arm with stated limits}
  \begin{columns}[c]
    \begin{column}{0.58\textwidth}
      \footnotesize
      \begin{tabular}{@{}lrl@{}}
        \toprule
        \textbf{Parameter} & \textbf{Value} & \textbf{Category} \\
        \midrule
        Loaded mass & 368 kg & derived \\
        Drive limit $v$ & 1.5 m/s & illustrative \\
        Deceleration $a_{\text{brake}}$ & 2 m/s$^2$ & illustrative \\
        Clear distance $D_{\text{clear}}$ & 1.10 m & illustrative \\
        Pre-brake delay $\tau_{\text{delay}}$ & 133.6 ms & derived \\
        Intent model & 336 ms & derived \\
        \bottomrule
      \end{tabular}\\[0.2cm]
      {\footnotesize Six of the card's fifteen values. A label says what a number can support.}
    \end{column}
    \begin{column}{0.37\textwidth}
      \centering\safeimg[width=\textwidth,height=4.2cm,keepaspectratio]{erflex-mobile-manipulator.jpg}\\
      {\footnotesize A real machine of this class}
    \end{column}
  \end{columns}
  \mlsyscredit{The Causal Boundary, \S1.4.2, table 2. Photo: ER-FLEX, Auledas, CC BY 4.0, Wikimedia Commons; its specifications differ from the card.}
  \note{[2.5 min] NARRATE: every lecture uses one machine, a warehouse mobile manipulator: a wheeled base, an arm, navigation lidar, and wrist cameras. It collects objects from a conveyor, opens a spring-latched cage door, and drives aisles shared with people. Its card fixes fifteen values; six are here. Read the category column: measured, derived, datasheet, illustrative, or chosen. Illustrative values show the arithmetic and support no safety margin; a chosen value is a design decision that must be checked against what it has to fit. The intent model, a 7-billion-parameter vision-language model, takes 336 ms per evaluation, so it runs at 2 Hz. ENGAGE: ask which values a real team would have to measure first. Expected: the deceleration and the clear distance, since every stop depends on them. WARN: the photo is a real machine of the same class; the card is the book's reference design and does not describe this product. FLEX: [CORE]. Lecture 1 owns the card; Lecture 3 owns the two-speed brain.}
\end{frame}

\begin{frame}{Does the base stop inside 1.10 m? Delay plus braking says yes}
  \begin{columns}[c]
    \begin{column}{0.42\textwidth}
      \begin{mlsyscard}{mlsysaccent}\footnotesize
        \textbf{Given}: $v = 1.5$ m/s, $a_{\text{brake}} = 2$ m/s$^2$,\\
        $D_{\text{clear}} = 1{,}100$ mm,\\ $\tau_{\text{delay}} = 133.6$ ms
      \end{mlsyscard}
    \end{column}
    \begin{column}{0.54\textwidth}
      \small
      \uncover<1->{\textcolor{delaycol}{Delay: $v\,\tau_{\text{delay}} = 1.5 \times 0.1336 = 200.4$ mm}}\\[0.05cm]
      \uncover<2->{\textcolor{plantcol}{Braking: $v^2/(2a_{\text{brake}}) = 2.25/4 = 562.5$ mm}}\\[0.05cm]
      \uncover<3->{\textbf{Stop: $d_{\text{stop}} = 762.9 \le 1{,}100$ mm}}
    \end{column}
  \end{columns}
  \vspace{0.15cm}
  \centering\mlsyssteps[width=0.95\textwidth,height=3.6cm,keepaspectratio]{stop_bar_step}{3}
  \mlsyscredit{The Causal Boundary, equation 1 and Napkin Math 1.1 (illustrative and derived inputs); computed in images/calcs.py.}
  \note{[3.5 min] LINK: the aisle from the opening. NARRATE: worked example. The elementary stopping budget has two terms: the distance traveled during the pre-brake delay $\tau_{\text{delay}}$, and the braking distance at the credible deceleration $a_{\text{brake}}$. The delay is 51.6 ms of frame age, a 60 ms chunk lease, a 1 ms permission tick, a 1 ms bus cycle, and 20 ms of brake onset: 133.6 ms. Step 1: 1.5 m/s times 0.1336 s is 200.4 mm, in amber. Step 2: 1.5 squared over 2 times 2 is 562.5 mm of braking. Step 3: the stop is 762.9 mm, inside the 1,100 mm clear distance with 337.1 mm to spare. ENGAGE: ask whether 337.1 mm is a safety margin. Expected: not yet; it ignores localization error, tracking error, a protective margin, and the longer stop of a jerk-limited brake. WARN: every input is illustrative or derived, so the answer is illustrative too. FLEX: [CORE]. Lecture 1 owns this budget; Lecture 12 adds the allowances and sets the aisle speed at 1.3 m/s.}
\end{frame}

\begin{frame}{Above 0.53 m/s, braking travel exceeds delay travel}
  \centering
  \safeimg[width=0.92\textwidth,height=5.4cm,keepaspectratio]{delay_vs_speed.pdf}
  \mlsyscredit{The Causal Boundary, \S1.4.3, figure 4; crossover $v^{*} = 2a_{\text{brake}}\tau_{\text{delay}}$, computed in images/calcs.py.}
  \note{[2 min] NARRATE: the two terms grow differently with speed. Delay travel, amber, grows linearly; braking travel grows with the square of speed; their sum is the black stop curve. Setting them equal gives the crossover speed, 2 times 2 times 0.1336, or 0.53 m/s. At the 1.5 m/s drive limit braking sets most of the stop (2.8 times the delay travel), and the stop is 762.9 mm, under the red clear-distance line. ENGAGE: ask which lever helps below the crossover. Expected: shortening the delay, since delay travel is the larger term there. WARN: dominance does not decide which change saves more; that depends on the reductions available. FLEX: [OPTIONAL]; if cut, state the crossover on the previous slide. Lecture 1 owns the crossover speed.}
\end{frame}

\begin{frame}{Estimate: what is the stop at half speed?}
  \mlsysactivity{Estimate}{3 min}
  \begin{columns}[c]
    \begin{column}{0.55\textwidth}
      \normalsize Same machine, same 133.6 ms delay, same 2 m/s$^2$.\\[0.2cm]
      Speed 0.75 m/s. What is $d_{\text{stop}}$?\\[0.3cm]
      \uncover<2->{\small Delay travel halves: 100.2 mm\\
      Braking quarters: 140.6 mm\\[0.1cm]
      {\normalsize\textbf{\textcolor{mlsysaccent}{Stop: 240.8 mm, about one third of 762.9}}}}
    \end{column}
    \begin{column}{0.42\textwidth}
      \centering\safeimg[width=\textwidth,height=4.2cm,keepaspectratio]{erflex-mobile-manipulator.jpg}
    \end{column}
  \end{columns}
  \mlsyscredit{The Causal Boundary, Napkin Math 1.1; computed in images/calcs.py. Photo: Auledas, CC BY 4.0, Wikimedia Commons.}
  \note{[3 min] NARRATE: two minutes, napkin only, pairs. Step 2: 0.75 times 0.1336 is 100.2 mm; 0.75 squared over 4 is 140.6 mm; 240.8 mm in all, 3.2 times shorter than at full speed. ENGAGE: ask a pair that got 381 mm what they did. Expected: they halved both terms. WARN: halving speed halves the delay term but quarters the braking term, which is why slowing down is the most powerful lever a designer holds. FLEX: [CORE].}
\end{frame}

\begin{frame}{The learned backbone is 22 ms of the 133.6 ms delay}
  \centering
  \safeimg[width=0.95\textwidth,height=3.6cm,keepaspectratio]{prebrake_bar.pdf}\\[0.3cm]
  {\small Each added millisecond anywhere in the chain costs 1.5 mm at the drive limit.}
  \mlsyscredit{The Causal Boundary, \S1.4.3 and Napkin Math 1.1; computed in images/calcs.py.}
  \note{[2 min] NARRATE: the 133.6 ms by part. The learned perception backbone, orange, is 22 ms, or 33 mm of travel. The rest of the frame's age is 29.6 ms. The chunk lease, green, is 60 ms, 90 mm: the longest an admitted chunk of motion may keep driving before the permission path must act again, which matters when the learned proposer stalls. Tick and bus are 2 ms; brake onset is 20 ms. ENGAGE: ask which design change shortens the delay most. Expected: the lease, or a separate protective path. A certified safety scanner wired to the drives bypasses perception and the lease entirely; industrial driverless trucks use one today (ISO 3691-4). WARN: a slower, more accurate model is not free; every millisecond it adds on this path is distance. FLEX: [CORE]. Lecture 4 owns the chunk lease.}
\end{frame}

\mlsysnote{[1 min] NARRATE: the sentence to carry from this segment. Delay is distance: the machine travels through every millisecond of sensing, computing, and signaling. Because that travel cannot be recalled, the path that can refuse an action must act before the clearance is spent, and it must not wait for the learned model. FLEX: [CORE]. Lecture 1 states this as the first law.}
\mlsysfocus{Delay is distance}{%
  $d_{\text{stop}} = v\,\tau_{\text{delay}} + \dfrac{v^2}{2a_{\text{brake}}} \;\le\; D_{\text{clear}}$\\[0.6cm]
  {\small\textcolor{white!75}{An independent path must be able to refuse an action before the clearance runs out.}}%
}

% =============================================================================
% SECTION: THE MACHINE
% =============================================================================
\mlsysnote{[0.5 min] NARRATE: stopping clearance is one margin. The same machine spends others. FLEX: [CORE].}
\mlsysdivider{Which margin does each task spend, and what is the machine made of?}{The Causal Boundary, \S1.5 and \S1.7}
\section{The Machine}

\begin{frame}{Each task spends a different margin: clearance, force, or heat}
  \begin{columns}[T]
    \begin{column}{0.32\textwidth}
      \centering\safeimg[width=\textwidth,height=2.6cm,keepaspectratio]{erflex-mobile-manipulator.jpg}\\[0.05cm]
      {\small\textbf{Mobility}}\\{\footnotesize stopping clearance}\\[0.05cm]
      {\footnotesize 762.9 of 1,100 mm}
    \end{column}
    \begin{column}{0.32\textwidth}
      \centering\safeimg[width=\textwidth,height=2.6cm,keepaspectratio]{franka-arm.jpg}\\[0.05cm]
      {\small\textbf{Contact}}\\{\footnotesize force on a latch}\\[0.05cm]
      {\footnotesize 39 N at brake command, 100 N limit}
    \end{column}
    \begin{column}{0.32\textwidth}
      \centering\safeimg[width=\textwidth,height=2.6cm,keepaspectratio]{bldc-stator-winding.jpg}\\[0.05cm]
      {\small\textbf{Holding}}\\{\footnotesize winding heat}\\[0.05cm]
      {\footnotesize 87 N$\cdot$m derated to 73.0}
    \end{column}
  \end{columns}
  \mlsyscredit{The Causal Boundary, \S1.5; Napkin Math 10.1 and 2.6 (illustrative). Photos: Auledas, CC BY 4.0; Ims, CC BY-SA 4.0; Medvedev, CC BY-SA 3.0; Wikimedia Commons.}
  \note{[2.5 min] NARRATE: the book names three machine classes by the margin each spends. Mobility spends stopping clearance: the base in the aisle, 762.9 of 1,100 mm. Contact spends force and deformation margin: when the arm presses a spring latch rated to 100 N, a 15 N tripwire and a 2 ms local contact loop brake the arm at 39 N; a stop that waited on the learned proposer would let the spring model reach 2,400 N. Holding spends heat and supply margin: an arm holding a load still does no mechanical work, yet the stall current heats the windings, and to protect their insulation a joint rated for 87 N$\cdot$m must be derated 16.1 percent, to 73.0 N$\cdot$m. The photos are a mobile manipulator, a Franka arm, and a motor's copper windings. ENGAGE: ask which class applies when the robot docks against a station. Expected: contact, briefly, then mobility again. WARN: these are regimes within one task; one robot passes through all three. FLEX: [CORE]. Lecture 1 owns the three classes.}
\end{frame}

\begin{frame}{Which margin runs out first in each task?}
  \mlsysactivity{Discuss}{2.5 min}
  \centering
  \begin{minipage}{0.86\textwidth}
    \large
    \choice{1}{The base drives toward a rack end at 1.5 m/s}\\[0.25cm]
    \choice{2}{The arm opens a spring-latched cage door}\\[0.25cm]
    \choice{3}{The arm holds a box still for one minute}\\[0.45cm]
    \normalsize\uncover<2->{Clearance, contact force, and heat. A machine with no motion can still run out of margin.}
  \end{minipage}
  \note{[2.5 min] NARRATE: two minutes with a neighbor; each pair names the margin for each task and one sensor that would watch it. Step 2: clearance, contact force, heat. Expected sensors: lidar and wheel encoders for clearance; a force-torque sensor or motor current for the latch; a winding temperature estimate for holding. ENGAGE: ask which task a model's accuracy affects least. Expected: holding still; the binding limit is thermal, and no prediction cools a winding. WARN: students assume a stationary arm is safe because nothing moves. FLEX: [OPTIONAL]; if cut, ask the third question aloud on the previous slide.}
\end{frame}

\begin{frame}{The machine has five levels and one line that only proposals cross}
  \begin{columns}[c]
    \begin{column}{0.58\textwidth}
      \centering
      \mlsyssteps[width=\textwidth,height=5.4cm,keepaspectratio]{\mlsyskit/stack_build}{7}
    \end{column}
    \begin{column}{0.39\textwidth}
      \small
      \only<1|handout:0>{\textbf{World}\\Contact, friction, payload, people. Continuous.}%
      \only<2|handout:0>{\textbf{Boundary}\\Register writes and sensor captures. Event-timed.}%
      \only<3|handout:0>{\textbf{Body}\\Drives, gears, windings, rails. Current loop 10 to 25 kHz.}%
      \only<4|handout:0>{\textbf{Nervous System}\\Permission path, lease, fallback. 1 kHz in one implementation.}%
      \only<5|handout:0>{\textbf{Brain}\\Learned proposers. Intent 2 Hz, chunks 20 Hz, setpoints 40 Hz.}%
      \only<6|handout:0>{\textbf{Proposal boundary}\\Above: may only propose.\\Below: may command.}%
      \only<7-|handout:1>{\textbf{Governance}\\An envelope around all five levels: who may command, on what evidence, until when.\\[0.2cm]
        \footnotesize Rates: World continuous; Body 10 to 25 kHz; Nervous System 1 kHz; Brain 2 to 40 Hz.}%
    \end{column}
  \end{columns}
  \note{[3.5 min] NARRATE: the course map; every lecture's dividers show it with its own level lit. Build bottom up. Step 1, World: the environment, which evolves whether or not the machine measures it. Step 2, Boundary: the transducers at both edges, register writes that energize windings and sensor captures that turn light and force into numbers; the actuation half is the causal boundary. Step 3, Body: drives, transmissions, windings, power rails; its current loops run at 10 to 25 kHz. Step 4, Nervous System: buses, schedulers, and the permission path, which checks each proposal and holds a fallback; 1 kHz in one implementation. Step 5, Brain: the learned proposers; the intent model at 2 Hz, chunk renewal at 20 Hz, setpoints at 40 Hz. Step 6: the proposal boundary, the line between a model that may only propose and a permission path that alone may let an action reach an actuator. Step 7: governance, the lifecycle envelope around all five levels; it is not a sixth level. ENGAGE: ask students to say what each boundary decides. Expected: the proposal boundary decides who may approve an action; the causal boundary is where an approved action becomes work. WARN: students merge the two boundaries; keep them apart in every answer. FLEX: [CORE]. Lecture 1 owns the five levels and the proposal boundary.}
\end{frame}

% =============================================================================
% SECTION: EVIDENCE
% =============================================================================
\mlsysnote{[0.5 min] NARRATE: a permission path needs limits to check against, and the machine needs evidence that it may operate at all. How much evidence is enough? FLEX: [CORE].}
\mlsysdivider{Why can a machine pass every test and still lack the evidence to operate?}{Closed-Loop Evaluation, \S7.6 and \S7.9; Supervisory Intervention, \S14.1}
\section{Evidence}

\begin{frame}{A decade of filings still shows 174 takeovers a year}
  \begin{columns}[c]
    \begin{column}{0.33\textwidth}
      \centering\safeimg[width=\textwidth,height=3.4cm,keepaspectratio]{waymo-jaguar-sf.jpg}\\
      {\footnotesize A Waymo vehicle in San Francisco, 2023, no one in the driver's seat}
    \end{column}
    \begin{column}{0.64\textwidth}
      \centering\safeimg[width=\textwidth,height=4.9cm,keepaspectratio]{dmv_waymo.pdf}
    \end{column}
  \end{columns}
  \mlsyscredit{Waymo annual filings, California DMV disengagement reports 2015 to 2025, via MLSysIM Literature.CaliforniaDisengagements; Supervisory Intervention, \S14.1, figure 1. Photo: Dllu, CC BY 4.0, Wikimedia Commons.}
  \note{[2.5 min] NARRATE: California requires every company testing autonomous vehicles on its public roads to report autonomous miles and disengagements, transfers of control to a person. These are Waymo's filings. With a safety driver, miles per disengagement rose from 1,244 in 2015 to 19,234 in 2025, yet the fleet still reported 174 takeovers in 2025, and the series is far from smooth. Driverless testing reports one disengagement every 71,815 miles, 54 in 2025. ENGAGE: ask what the 2021 drop means. Expected: it may reflect where the fleet drove and what it chose to count; the rate is not a capability curve. WARN: each company defines which disengagements are reportable; the regulator warns against comparing companies. FLEX: [CORE]. Human authority is a permanent part of the design; Lecture 14 owns takeover delay.}
\end{frame}

\begin{frame}{Estimate: how many clean runs support 99 percent?}
  \mlsysactivity{Estimate}{2.5 min}
  \centering
  \normalsize The arm opens the cage door 20 times with no failure.\\
  That supports a success rate of at least 86.1\% at 95\% confidence.\\[0.4cm]
  How many clean runs would support 99\%?\\[0.3cm]
  \begin{minipage}{0.5\textwidth}
    \Large
    \choice{A}{30}\\[0.1cm]
    \choice{B}{100}\\[0.1cm]
    \alt<2->{\colorbox{mlsysaccent!15}{\choice{C}{about 300}}}{\choice{C}{about 300}}\\[0.1cm]
    \choice{D}{about 3,000}
  \end{minipage}
  \mlsyscredit{Closed-Loop Evaluation, Napkin Math 7.1 and \S7.6; computed in images/calcs.py.}
  \note{[2.5 min] NARRATE: one minute with a neighbor. Hint: the chance of $n$ clean runs if the true success rate is $p$ is $p^n$; the bound is the $p$ at which that chance falls to 5 percent. Step 2: C, 299 runs. Twenty clean runs give $p^{20} = 0.05$, so $p = 0.861$. For 99 percent, $0.99^n = 0.05$ gives $n = 299$. ENGAGE: ask a student who chose D what they computed. Expected: the count for 99.9 percent, which is next. WARN: students read 20 of 20 as 100 percent; a clean record bounds the rate, it does not measure it. FLEX: [CORE].}
\end{frame}

\begin{frame}{Each added nine of reliability costs ten times the clean runs}
  \centering
  \mlsyssteps[width=0.92\textwidth,height=4.3cm,keepaspectratio]{clean_runs_step}{4}\\[0.15cm]
  \uncover<4->{\small At 40 s per trial: 20 runs take 13.3 min; 2,995 runs take 33.3 h of cell time.}
  \mlsyscredit{Closed-Loop Evaluation, Napkin Math 7.1 and 7.2 (40 s per trial, illustrative); computed in images/calcs.py.}
  \note{[2.5 min] NARRATE: worked example 2. Each bar is the number of consecutive clean runs needed for a lower success limit at 95 percent confidence, on a log axis. Step 1: 20 runs, 86.1 percent. Step 2: 29 runs for 90 percent. Step 3: 299 for 99 percent. Step 4: 2,995 for 99.9 percent. At 25 s of motion and 15 s of reset per trial, the first twenty take 13.3 minutes and 2,995 take 33.3 hours of cell time, before any cooldown for the windings. ENGAGE: ask what a hazard rate of one in a million hours would need. Expected: about 3 million failure-free hours, beyond any test campaign. WARN: the bound assumes independent, identical trials; a campaign this long outlives that assumption. FLEX: [CORE]. Lecture 7 owns the zero-failure bound.}
\end{frame}

\begin{frame}{Driving cannot show fatality rates; 56.7 million miles is short}
  \centering
  \safeimg[width=0.95\textwidth,height=4.4cm,keepaspectratio]{rand_miles.pdf}\\[0.2cm]
  {\small The fatality bar is 275 million miles, 4.9 times Waymo's rider-only total.}
  \mlsyscredit{Kalra and Paddock, Driving to Safety, RAND 2016; Kusano et al. 2025 (rider-only miles through January 2025); Closed-Loop Evaluation, \S7.9.1, figure 5.}
  \note{[1.5 min] NARRATE: the same arithmetic at fleet scale. At 1.09 fatalities per 100 million miles, a fleet must log 275 million failure-free miles to bound its fatality rate at the human level with 95 percent confidence. Waymo's 56.7 million rider-only miles through January 2025 clear the bars for injuries and crashes, and published comparisons show large reductions in injury crashes, but the miles fall short of the fatality bar by a factor of 4.9. ENGAGE: ask what has to supply the rest of the evidence. Expected: runtime limits enforced by an independent path, fault injection, and an argument for each claim. FLEX: [OPTIONAL]. Lecture 7 owns the exposure wall.}
\end{frame}

\begin{frame}{Four laws organize the course, one per part}
  \centering\small Only the first is a physical law. The other three are design principles the book raises to laws.\par\vspace{0.25cm}
  \footnotesize
  \begin{columns}[T]
    \begin{column}{0.24\textwidth}
      \begin{mlsyscard}{conseq}
        \textbf{1. Irreversibility}\\[0.05cm]Software cannot roll back physical consequences.\\[0.1cm]
        \textcolor{midgray}{Physical law; Part I}
      \end{mlsyscard}
    \end{column}
    \begin{column}{0.24\textwidth}
      \begin{mlsyscard}{datastroke}
        \textbf{2. Endogenous data}\\[0.05cm]A machine's actions choose its next observations.\\[0.1cm]
        \textcolor{midgray}{Covariate shift; Part II}
      \end{mlsyscard}
    \end{column}
    \begin{column}{0.24\textwidth}
      \begin{mlsyscard}{nervecol}
        \textbf{3. Proposal is not permission}\\[0.05cm]Only an independent path may admit a proposal.\\[0.1cm]
        \textcolor{midgray}{Runtime assurance; Part III}
      \end{mlsyscard}
    \end{column}
    \begin{column}{0.24\textwidth}
      \begin{mlsyscard}{brainnavy}
        \textbf{4. Permission must be earned}\\[0.05cm]Authority holds per condition, on evidence, and expires.\\[0.1cm]
        \textcolor{midgray}{Assurance case; Part IV}
      \end{mlsyscard}
    \end{column}
  \end{columns}
  \mlsyscredit{The Causal Boundary, \S1.10 and table 3.}
  \note{[2.5 min] NARRATE: the course in four sentences. The first law, irreversibility: work delivered across the causal boundary cannot be recalled, so every check must finish before the clearance is spent; that is today's stopping budget. The second, endogenous data: a robot's actions choose what it sees next, so a policy's competence holds only over the states its data and closed-loop tests covered; its established name is covariate shift (Ross, Gordon, and Bagnell 2011). The third, proposal is not permission: a learned model may propose motion, and only an independent permission path may admit it; its name is runtime assurance, from the Simplex architecture (Sha 2001). The fourth, permission must be earned: authority to operate is granted per operating condition on recorded evidence and lapses with it; its name is the assurance case (UL 4600). ENGAGE: ask which law today's clean-run ladder serves. Expected: the fourth. WARN: never call the second to fourth laws physical; the book raises them to laws by choice because breaking them carries predictable risk. FLEX: [CORE]. Lecture 1 teaches the four laws; each part's first lecture states its law.}
\end{frame}

% =============================================================================
% SECTION: THE COURSE
% =============================================================================
\mlsysnote{[0.5 min] NARRATE: back to the opening question, then the term ahead. FLEX: [CORE].}
\mlsysdivider{How far does a robot travel while it decides, and what comes next?}{The Causal Boundary, \S1.11; Reader Guide}
\section{The Course}

\begin{frame}{Close your notes: what does the aisle now show?}
  \mlsysactivity{Retrieval}{2.5 min}
  \begin{columns}[c]
    \begin{column}{0.58\textwidth}
      \centering\safeimg[width=\textwidth,height=3.6cm,keepaspectratio]{aisle_step3.pdf}
    \end{column}
    \begin{column}{0.39\textwidth}
      \small Write down:\\[0.1cm]
      \choice{a}{the stopping budget}\\
      \uncover<2->{\textcolor{mlsysaccent}{$v\tau_{\text{delay}} + v^2/2a_{\text{brake}} \le D_{\text{clear}}$}}\\[0.12cm]
      \choice{b}{the two boundaries}\\
      \uncover<2->{\textcolor{mlsysaccent}{who may approve; where work begins}}\\[0.12cm]
      \choice{c}{the four laws}\\
      \uncover<2->{\textcolor{mlsysaccent}{irreversibility, endogenous data, proposal is not permission, permission must be earned}}
    \end{column}
  \end{columns}
  \mlsyscredit{The Causal Boundary, Napkin Math 1.1; Sensor Perception, Napkin Math 8.2; computed in images/calcs.py.}
  \note{[2.5 min] NARRATE: 90 seconds, notes closed, written answers; then Step 2. Return to the opening question with the aisle complete. One frame costs 77.4 mm. The whole elementary stop is 200.4 mm of delay plus 562.5 mm of braking, 762.9 mm inside the 1,100 mm clear distance. (b) The proposal boundary decides who may approve an action; the causal boundary is where an approved action becomes work. (c) The four laws in order. ENGAGE: collect one answer for (b) and ask which level each boundary sits beside. Expected: the proposal boundary between Brain and Nervous System; the causal boundary at the Boundary level. FLEX: [CORE]. Lecture 12 turns this spare into the chosen aisle speed of 1.3 m/s.}
\end{frame}

\begin{frame}{Eighteen lectures follow the book's four parts}
  \centering
  \mlsyskitfig[width=\textwidth,height=6.0cm,keepaspectratio]{course_order}
  \note{[1 min] NARRATE: the term in order, one card per lecture, grouped by part, each card naming its level. Lecture 1 opens with the causal boundary. Part I measures the body, the brain, and the nervous system that joins them. Part II teaches the machine: data, training, and closed-loop evaluation. Part III runs it: perception through placement. Part IV governs it: intervention, verification, release, and responsibility. Lecture 18 closes. ENGAGE: ask why the body comes before the brain. Expected: the limits a proposal must respect are set by the body, so they come first. FLEX: [OPTIONAL].}
\end{frame}

\begin{frame}{Each lecture adds one record to a chain of evidence}
  \centering
  \safeimg[width=0.95\textwidth,height=5.6cm,keepaspectratio]{record_chain.pdf}
  \mlsyscredit{After The Causal Boundary, \S1.11, figure 9.}
  \note{[1 min] NARRATE: every lecture from 2 to 18 adds one record, a checkable artifact that later decisions rely on. The limit record states what the body can do, under what conditions, on what evidence; the handoff record states the lease and tick. Part II's records state what the proposer has shown; Part III's records carry live observations and proposals to the enforcement record, where the two strands meet; Part IV's records grant and limit authority. Lecture 3 adds none; its limits enter the handoff record. Lecture 18 keeps a register of the claims still open. ENGAGE: ask which record today's 762.9 mm would start. Expected: a limit record, still missing its uncertainty, detector, and measured provenance. FLEX: [OPTIONAL]. Lecture 1 owns the record chain.}
\end{frame}

\begin{frame}{By the end of the course, you will be able to}
  \begin{columns}[T]
    \begin{column}{0.56\textwidth}
      \small
      \begin{itemize}\setlength{\itemsep}{0.02cm}
        \item say what response remains if a proposal is wrong, stale, or missing;
        \item size the stopping budget, chunk lease, and frame age;
        \item state what data and tests do and do not support;
        \item design a path that admits, corrects, or refuses each proposal;
        \item write a release argument with named owners and an expiry.
      \end{itemize}
    \end{column}
    \begin{column}{0.20\textwidth}
      \centering\safeimg[width=\textwidth,height=4.0cm,keepaspectratio]{tiago-mobile-manipulator.jpg}
    \end{column}
    \begin{column}{0.21\textwidth}
      \begin{mlsyscard}{mlsysaccent}\scriptsize
        \textbf{You need}: Python, ML intuition, $F = ma$, basic probability.
      \end{mlsyscard}\par\vspace{0.05cm}
      \begin{mlsyscard}{BlueLine}\scriptsize
        \textbf{Three paths in}: ML, embedded systems, robotics and control.
      \end{mlsyscard}\par
      {\scriptsize\textcolor{midgray}{Schedule and policies: the course page.}}
    \end{column}
  \end{columns}
  \mlsyscredit{Photo: TIAGo++ mobile manipulator, PAL Robotics, CC BY-SA 4.0, Wikimedia Commons. Outcomes: Volume IV course plan; prerequisites: the Reader Guide.}
  \note{[1.5 min] NARRATE: five outcomes, one per part plus the opening: the response if the proposal fails (Part I and today), the budgets (Part I), evidence from data and tests (Part II), the permission path (Part III), and the release argument (Part IV). Prerequisites are the book's: programming, machine learning intuition, freshman mechanics, and elementary probability; no control theory, fleet experience, or other volume is assumed. The photo is TIAGo++, a two-armed mobile manipulator used in research labs. Students arrive by three paths. Machine learning students will find back-EMF, real-time scheduling, and barrier functions hard; embedded students, imitation learning and diffusion decoders; robotics students, confidence bounds. The appendices on machine learning, control, and systems cover each. FLEX: [CORE]. Replace the last line with your schedule, grading, and labs.}
\end{frame}

\begin{frame}{Takeaways}
  \begin{columns}[c]
    \begin{column}{0.68\textwidth}
      \small
      \begin{itemize}\setlength{\itemsep}{0.02cm}
        \item One camera frame costs 77.4 mm of travel at 1.5 m/s.
        \item Physical AI joins learning, physical feedback, and delegated authority.
        \item The elementary stop is 762.9 mm of 1,100; each millisecond adds 1.5 mm.
        \item Five levels; the proposal boundary separates proposing from permitting.
        \item A 99.9\% bound needs 2,995 clean runs; tests alone cannot earn permission.
        \item Four laws, one per part; one record per lecture.
      \end{itemize}
    \end{column}
    \begin{column}{0.29\textwidth}
      \centering
      \safeimg[width=\textwidth,height=0.78cm,keepaspectratio]{thumb_aisle_step2.pdf}\\[0.04cm]
      \safeimg[width=\textwidth,height=0.78cm,keepaspectratio]{thumb_epoch_compute_step2.pdf}\\[0.04cm]
      \safeimg[width=\textwidth,height=0.78cm,keepaspectratio]{thumb_stop_bar_step3.pdf}\\[0.04cm]
      \safeimg[width=\textwidth,height=0.78cm,keepaspectratio]{thumb_stack.pdf}\\[0.04cm]
      \safeimg[width=\textwidth,height=0.78cm,keepaspectratio]{thumb_clean_runs_step4.pdf}\\[0.04cm]
      \safeimg[width=\textwidth,height=0.78cm,keepaspectratio]{thumb_record_chain.pdf}
    \end{column}
  \end{columns}
  \note{[1.5 min] NARRATE: one line per segment, each beside its picture. Read them as the answer to the driving question: a robot travels 77.4 mm per frame and 762.9 mm per elementary stop at 1.5 m/s, so an independent path must be able to refuse in time, and the evidence that it may operate grows expensive fast. FLEX: [CORE].}
\end{frame}

\begin{frame}{Open questions to carry through the course}
  \begin{columns}[c]
  \begin{column}{0.66\textwidth}
  \small
  \begin{itemize}\setlength{\itemsep}{0.3cm}
    \item If a model is right 999 times in 1,000, what should the machine do on the thousandth, and who decides?
    \item Which will limit physical AI first: computation on the robot, data from robots, or evidence about robots?
    \item When a robot and a person share an aisle, should the robot's uncertainty cost the robot speed or the person attention?
  \end{itemize}
  \end{column}
  \begin{column}{0.31\textwidth}
    \centering\safeimg[width=\textwidth,height=3.6cm,keepaspectratio]{waymo-jaguar-sf.jpg}
  \end{column}
  \end{columns}
  \mlsyscredit{Photo: a Waymo vehicle with no one in the driver's seat, San Francisco, 2023, Dllu, CC BY 4.0, Wikimedia Commons.}
  \note{[1 min] NARRATE: no single answer. Ask each student to write one sentence on any of the three and keep it; Lecture 18 returns to the course's opening question with every record assembled. FLEX: [OPTIONAL].}
\end{frame}

\begin{frame}{Read and explore}
  \begin{columns}[T]
    \begin{column}{0.56\textwidth}
      \footnotesize
      \begin{itemize}\setlength{\itemsep}{0.08cm}
        \item \textbf{Read:} The Causal Boundary, \S1.1 to \S1.11; the Reader Guide.
        \item Sha, ``Using Simplicity to Control Complexity,'' \emph{IEEE Software}, 2001.
        \item Ross, Gordon, and Bagnell, ``A Reduction of Imitation Learning \ldots,'' AISTATS 2011.
        \item Kalra and Paddock, \emph{Driving to Safety}, RAND, 2016.
        \item UL Standards \& Engagement, ANSI/UL 4600, 2023.
        \item NTSB, HAR-19/03, Tempe collision report, 2019.
      \end{itemize}
    \end{column}
    \begin{column}{0.41\textwidth}
      {\footnotesize\textbf{Explore}: the public filings behind today's plot (Waymo rows)}\par\vspace{0.1cm}
      \scriptsize
      \begin{tabular}{@{}rrrr@{}}
        \toprule
        Year & Miles & Diseng. & Mi/diseng. \\
        \midrule
        2015 & 424,331 & 341 & 1,244 \\
        2020 & 628,839 & 21 & 29,945 \\
        2021 & 2,325,843 & 292 & 7,965 \\
        2025 & 3,346,709 & 174 & 19,234 \\
        \bottomrule
      \end{tabular}\\[0.05cm]
      {\scriptsize Safety-driver permit only.}
    \end{column}
  \end{columns}
  \mlsyscredit{California DMV autonomous vehicle disengagement reports (dmv.ca.gov), Waymo annual filings, as stored in MLSysIM Literature.CaliforniaDisengagements; 2015 transcribed from the filed PDF.}
  \note{[1 min] NARRATE: the chapter first; it develops every idea from today. Sha's paper is the root of the third law; Ross, Gordon, and Bagnell of the second; RAND and UL 4600 of the fourth. The NTSB report on the 2018 Tempe collision opens Lecture 1; reading it first pays off. The table is four rows of the public filings behind today's plot: students can download the reports from the California DMV and recompute miles per disengagement themselves. The 2020 to 2021 swing, 29,945 to 7,965 while miles grew 3.7 times, shows why a reported rate is not a capability measure. FLEX: [CORE].}
\end{frame}

\begin{frame}[c]
  \begin{columns}[c]
    \begin{column}{0.66\textwidth}
      {\small\textcolor{lightgray}{NEXT LECTURE \textbar{} THE CAUSAL BOUNDARY}}\\[0.5cm]
      {\Large\bfseries Why can a better prediction still produce a worse physical action?}
    \end{column}
    \begin{column}{0.30\textwidth}
      \centering\safeimg[width=\textwidth,height=2.6cm,keepaspectratio]{thumb_stop_bar_step3.pdf}
    \end{column}
  \end{columns}
  \note{[0.5 min] NARRATE: close on the question Lecture 1 answers. It recalls today's budget in one slide and spends its time on the crossover speed, recovery, and a real collision. Ask students to read The Causal Boundary before then. FLEX: [CORE].}
\end{frame}

\end{document}
